% Context from chapters-tex/wavelets.tex; for mathematical reference, not compilation.
 same detail spaces).

\begin{prop}
	For a real filter $h\in\ell^1(\ZZ)$ satisfying~\eqref{eq-cond-h-2}, the filter $g \in \ell^2(\ZZ)$ defined by
	\eql{\label{eq-qmf}
		\foralls n \in \ZZ, \quad g_n = (-1)^{1-n}h_{1-n}
	}
	satisfies conditions~\eqref{eq-cond-g-1} and~\eqref{eq-cond-g-2}.
\end{prop}

\begin{proof}
	Taking the Fourier series of the defining relation gives
	\eql{\label{eq-qmf-fourier}
		\hat g(\om) = e^{-i \om} \hat h(\om+\pi)^*, 
	}
	so that 
	\eq{
		|\hat g(\om)|^2 + |\hat g(\om+\pi)|^2 = 
		|e^{-i \om} \hat h(\om+\pi)^*|^2 + |e^{-i (\om+\pi)} \hat h(\om+2\pi)^*|^2
		= 
		|\hat h(\om+\pi)|^2 + |\hat h(\om+2\pi)|^2 = 2.
	}
	where we used the fact that $\hat h$ is $2\pi$-periodic, 
	and also 
	\begin{align*}
		\hat g(\om) \hat h(\om)^* + \hat g(\om+\pi) \hat h(\om+\pi)^* &= 
		e^{-i \om} \hat h(\om+\pi)^* \hat h(\om)^* + e^{-i (\om+\pi)} \hat h(\om+2\pi)^* \hat h(\om+\pi)^* \\
		&=
		(e^{-i \om}+e^{-i (\om+\pi)}) \hat h(\om+\pi)^* \hat h(\om)^*
		=0.
	\end{align*}
\end{proof}





                                   
\subsection{Wavelet Design Constraints}

One can construct a multiresolution wavelet pair by designing a finite filter $h$ satisfying~\eqref{eq-cond-h-1}--\eqref{eq-cond-h-star}. The scaling function is obtained through the infinite product~\eqref{eq-dfn-cascade-inf}. Except in special cases such as Haar wavelets, it has no elementary closed form. The QMF relation~\eqref{eq-qmf} then defines $g$, and~\eqref{eq-rel-psi-phi} defines $\psi$.

Fourier atoms are fixed exponentials, whereas mother wavelets can be designed to meet several requirements:
\begin{rs}
	\item Size of the support.
	\item Cancellation of polynomials, measured by the number $p$ of vanishing moments.
	\item Symmetry (for compactly supported scalar orthogonal wavelets, the Haar system is the exceptional symmetric case).
	\item Smoothness (number of derivatives).
\end{rs}
We now explain how these design requirements interact with conditions~\eqref{eq-cond-h-1}--\eqref{eq-cond-g-2}.

  
\paragraph{Vanishing moments.}

A wavelet $\psi$ has $p$ vanishing moments if
\eql{\label{eq-dfn-vm}
	\foralls k=0,\ldots,p-1, \quad \int_\RR \psi(x) \, x^k \d x = 0.
}
Vanishing moments make $\dotp{f}{\psi_{j,n}}$ small when $f$ has regularity $\Cal$ with $\al<p$ on Supp$(\psi_{j,n})$. A Taylor expansion of $f$ about $2^j n$ on the wavelet's support explains this cancellation.
 
The vanishing-moment condition has the following equivalent Fourier characterization; see Figure~\ref{fig-vanmoments}.

\begin{figure}[tbp]
\centering
\BookDrawing{tikz/wavelets-vanishing-moments}
\caption{Vanishing moments correspond to vanishing derivatives of the Fourier transform.}
\label{fig-vanmoments}
\end{figure}

\begin{prop}
Assuming the required moments of $\psi$ and derivatives of the filter symbols exist, a wavelet obtained by the QMF construction~\eqref{eq-qmf} has $p$ vanishing moments if and only if
\eql{\label{eq-cond-vm-fourier}
	\foralls k=0,\ldots,p-1, \quad
	\frac{\d^k \hat h}{\d \om^k}(\pi) = \frac{\d^k \hat g}{\d \om^k}(0) = 0.
}
\end{prop}
\begin{proof}
	Integrability of the moments allows differentiation under the Fourier integral: $\hat\psi^{(k)} = \Ff( (-\imath \cdot)^k \psi(\cdot) )$. Thus
	\eq{
		(-\imath)^k \int_\RR x^k \psi(x) \d x  = \hat\psi^{(k)}(0).
	}
	Relations~\eqref{eq-rel-psi-phi} and~\eqref{eq-qmf} imply
	\eq{
		\hat \psi(2\om) = \hat h(\om+\pi)^* \rho(\om)
		\qwhereq
		\rho(\om) \eqdef \frac{1}{\sqrt{2}} e^{-i \om}  \hat \phi(\om).
	}
	and differentiating this relation shows	
	\eq{
		2 \hat \psi^{(1)}(2\om) = \hat h(\om+\pi)^* \rho^{(1)}(\om) + \hat h^{(1)}(\om+\pi)^* \rho(\om)
	}
	Since $\hat h(\pi)=0$ and $\rho(0)=\hat\phi(0)/\sqrt2\neq0$, this shows that $\hat\psi^{(1)}(0)=0 \Leftrightarrow \hat h^{(1)}(\pi)^*=0$.
	 
	Induction gives the higher-order equivalence
	$\hat\psi^{(k)}(0)=0 \Leftrightarrow \hat h^{(k)}(\pi)^*=0$, provided all derivatives of lower order vanish.
\end{proof}

Conditions~\eqref{eq-cond-h-1} and~\eqref{eq-cond-h-2} give $\hat h(\pi)=0$, so every admissible wavelet has at least $1$ vanishing moment: $\int \psi=0$. By~\eqref{eq-cond-vm-fourier}, additional vanishing moments require the Fourier transform $\hat h$ to vanish to higher order at $\om=\pi$.

  
\paragraph{Support.}

Figure~\ref{fig-vanishing-moments-1d} shows the wavelet coefficients of a piecewise smooth signal. The cancellation provided by $\psi$ suppresses coefficients in smooth regions, leaving the largest coefficients near singularities.

\myfigure{
\image{wavelets}{.6}{vanishing-moments-1d-coefs}
}{ 
	Location of large wavelet coefficients.  
}{fig-vanishing-moments-1d}

Choosing $\psi$ with short support limits how many wavelets intersect each singularity.
 
One can show that the size of the support of $\psi$ is proportional to the size of the support of $h$.
 
Short support competes with the vanishing-moment requirement~\eqref{eq-dfn-vm}. Indeed, one can prove that
for an orthogonal wavelet basis with $p$ vanishing moments
\eq{
	|\text{Supp}(\psi)| \geq 2p-1,
} 
where the left-hand side denotes the length of the smallest closed interval containing the support of a compactly supported orthogonal mother wavelet.


Chapter~\ref{chap-approx} studies in detail the tradeoff between support size and vanishing moments in nonlinear approximation of piecewise smooth signals.

  
\paragraph{Smoothness.}

In compression or denoising applications, an approximate signal is recovered from a partial set $I_M$ of coefficients,
\eq{
	f_M = \sum_{(j,n) \in I_M} \dotp{f}{\psi_{j,n}} \psi_{j,n}.
}
This finite sum is at least as smooth as $\psi$.

A smooth wavelet $\psi$ helps avoid visible reconstruction artifacts. Smoothness controls the regularity of the reconstructed signal, while vanishing moments and localization primarily determine the approximation rates studied below. These properties are often linked: in many families, including the Daubechies wavelets introduced next, more vanishing moments also bring greater smoothness.


                                   
\subsection{Daubechies Wavelets}

To build a wavelet $\psi$ with a fixed number $p$ of vanishing moments, one designs the filter $h$, and uses the quadrature mirror filter relation~\eqref{eq-qmf} to compute $g$. One therefore looks for $h$ such that
\eq{
	|\hat h(\om)|^2 + |\hat h(\om+\pi)|^2 = 2, \quad 
	\hat h(0)=\sqrt{2}, \qandq
	\foralls k=0,\ldots,p-1, \; {\frac{\d^k \hat h}{\d \om^k} }(\pi)=0.
}
These conditions give algebraic relations between the coefficients of $h$, which can be solved explicitly using the Euclidean division algorithm for polynomials.

The resulting Daubechies